\section{Methods}

\textit{Source of truth for numbers is \texttt{paper\_results.md} / \texttt{RESEARCH\_LOG.md} \S{}B. This document specifies the shared representation, encoder, pretraining objectives, transfer protocol, preprocessing, and the analyses (statistics, calibration, representation similarity, subject-scaling, controller). Written 2026-07-05, reflecting the v2 preprocessing and the multi-seed + Phase 1--3 analyses.}

\subsection{1. Overview of the study design}

We hold a single controlled variable --- the \textbf{pretraining objective} used to initialize a motion encoder --- fixed everything else, and measure what transfers from an abundant vision-captured skeleton source (NTU RGB+D) to a data-scarce body-worn inertial target (Xsens). Both modalities are mapped to one representation (Local Relative Quaternions, \S{}2) so that a single encoder architecture (\S{}3) serves both. We compare five initializations --- \texttt{scratch}, \texttt{mae}, \texttt{supMAE}, \texttt{supLP120}, and (as a scoped exhibit) \texttt{SupCon} (\S{}4) --- under a deployment-realistic subject-held-out few-shot calibration protocol (\S{}5), across a preprocessing that closes the modality gap (\S{}6). We then analyze the results along five axes: paired multi-seed statistics (\S{}7.1), clip-level significance via McNemar (\S{}7.2), calibration (\S{}7.3), representation similarity via CKA (\S{}7.4), subject-count scaling (\S{}7.5), and a downstream control-reliability simulation (\S{}8).

\subsection{2. Shared representation --- Local Relative Quaternions (LRQ)}

Every clip, from either modality, is encoded as a tensor of shape \texttt{(T, 17, 4)}: \texttt{T} temporal frames at 30 Hz, 17 body segments, and a unit quaternion per segment. Each segment's orientation is expressed \textbf{relative to its parent} in the kinematic hierarchy rather than in world space, making the representation invariant to global body rotation and translation.

\begin{itemize}
  \item \textbf{17-segment body model:} Pelvis, T8, Head, \{R,L\}-Shoulder, \{R,L\}-UpperArm, \{R,L\}-Forearm, \{R,L\}-Hand, \{R,L\}-UpperLeg, \{R,L\}-LowerLeg, \{R,L\}-Foot. Hierarchy: Pelvis$\rightarrow$T8$\rightarrow$Head; T8$\rightarrow$Shoulder$\rightarrow$UpperArm$\rightarrow$Forearm$\rightarrow$Hand; Pelvis$\rightarrow$UpperLeg$\rightarrow$LowerLeg$\rightarrow$Foot. The NTU joint topology is remapped to match the Xsens segment set exactly.
  \item \textbf{Flattening:} the \texttt{(17, 4)} per-frame orientation is flattened to a 68-d vector, so a clip is \texttt{(T, 68)}.
  \item \textbf{Temporal padding:} clips are zero-padded to \texttt{max\_frames = 120}; a boolean mask (1 = real frame, 0 = pad) travels with each clip and is converted to a Transformer padding mask so pooling (\texttt{masked\_mean}) averages only over real frames.
\end{itemize}

NTU quaternions are computed as bone rotations relative to an N-pose reference, projected through the hierarchy (\texttt{src/data/ntu\_parser.py}). Xsens quaternions are computed by the v2 pipeline (\S{}6).

\subsection{3. Encoder architecture}

The primary encoder (\texttt{KinematicEncoder}, \texttt{src/models/kinematic\_encoder.py}) is a Transformer over the temporal axis: Linear 68$\rightarrow$512, a learnable positional embedding \texttt{(1, 120, 512)}, and a 3-layer \texttt{TransformerEncoder} (pre-norm, 8 heads, 4$\times$ feed-forward). With \texttt{pool=False} it returns per-frame features \texttt{(B, T, 512)} for the reconstruction decoder; with \texttt{pool=True} it returns an L2-normalized clip embedding \texttt{(B, 512)} for classification. A padding mask \texttt{(B, T)} is threaded throughout.

The MAE decoder (\texttt{KinematicDecoder}) is a 1-layer Transformer + linear projection back to 68-d, used only during reconstruction pretraining. (A heavier MotionBERT-style \texttt{DSTformerQuatEncoder} exists in the codebase and was used in early experiments; all reported results use \texttt{KinematicEncoder}.)

\subsection{4. Pretraining objectives (the independent variable)}

All source pretraining is on NTU. Each objective produces one encoder checkpoint that seeds the transfer protocol (\S{}5).

\begin{center}
\small
\begin{tabular}{lll}
\toprule
id & objective & training signal \\
\midrule
\texttt{scratch} & none (control) & random init; no source pretraining \\
\texttt{mae} & masked autoencoding (pure reconstruction) & mask \textasciitilde{}70\% of frames, reconstruct the full quaternion sequence (geodesic quaternion loss `1 $-$ \\
\texttt{supMAE} & hybrid reconstruction + supervision & shared encoder, multi-task MAE + cross-entropy on NTU labels, CE warm-up and EMA loss balancing \\
\texttt{supLP120} & pure supervised & cross-entropy over the 120 NTU action classes \\
\texttt{SupCon} & supervised contrastive (Khosla et al. 2020) & contrastive loss with label-defined positives on NTU \\
\bottomrule
\end{tabular}
\end{center}

\texttt{supMAE} follows the vision-domain SupMAE formulation and is cited/renamed accordingly ("Sup+MAE (hybrid)") in prose. \texttt{SupCon} is included as a \textbf{scoped negative exhibit}, not a first-class arm: it was evaluated only under the earlier single-seed local preprocessing, where it underperformed \texttt{scratch} at every k>0; it is not part of the v2 multi-seed sweep and is reported as such (see Discussion). \texttt{DANN} gradient-reversal variants exist in the codebase (mid-pack under the earlier preprocessing) and are cited as a domain-adaptation reference point, not a headline arm.

\subsection{5. Transfer protocol --- subject-held-out k-shot calibration (LOSO)}

The Xsens target has 5 subjects (sub7--sub11) and 22 gestures. Each subject in turn is the fully held-out test subject (LOSO). For each held-out subject and each objective:

\begin{enumerate}
  \item \textbf{Base phase.} Initialize the encoder from the objective's checkpoint (or random for \texttt{scratch}), then fine-tune the full model on the 4 non-held-out subjects (80 epochs).
  \item \textbf{Calibration phase.} Freeze the encoder; take \texttt{k} labeled examples per class \textit{from the held-out subject}, train the classification head only, and evaluate on the remainder of that subject's clips. \texttt{k $\in$ \{0, 1, 3, 5, 10\}}; \texttt{k=0} = no calibration (zero-shot to the new subject after base training).
\end{enumerate}

This isolates the prior: \texttt{scratch}, \texttt{supMAE}, and \texttt{supLP120} are fine-tuned on the \textit{same} 4 subjects and calibrated identically, differing \textbf{only} in the initialization. Thus \texttt{supMAE $-$ scratch} and \texttt{supLP120 $-$ scratch} are apples-to-apples answers to the deployment question --- "does adding a cross-domain prior on top of the 4 available subjects adapt better to the 5th?"

\textbf{Discriminating regime.} All methods reach 94--98\% at k$\ge$5 (ceiling); k$\in$\{0,1,3\} is where objectives separate, so those are the primary reporting ks.

\subsection{6. Preprocessing --- v2 position-reconstructed Xsens}

The Xsens$\rightarrow$LRQ conversion is the second controlled lever. The paper's primary preprocessing (\textbf{v2}) rebuilds each Xsens segment orientation from the mvnx \textbf{position} streams through NTU's own shortest-arc \texttt{get\_bone\_quaternion} construction, plus a per-session T-pose world alignment (\texttt{src/scripts/IMU\_batch\_processor\_v2.py}, 2736 clips $\rightarrow$ \texttt{Data\_Processed/imu\_quats\_v2/}). Because it reconstructs orientation the way NTU constructs it (twist excluded identically), it closes the source$\rightarrow$target feature gap: MMD\textsuperscript{2} is \textit{below} the earlier measured-orientation baseline for every encoder.

Two earlier variants are retained only as gap reference points: \textbf{local} (v1: measured parent-relative orientation, including axial twist) and \textbf{swing} (axial twist stripped from Xsens post hoc). The swing variant was originally intended to \textit{close} the gap but \textit{widened} it \textasciitilde{}3$\times$ --- a symmetric test (twist-stripping NTU too) showed the increase came from twist removal damaging the Xsens signal, not from asymmetry. Swing is therefore reported as a \textbf{methodological finding} ("surface preprocessing does not do what intuition predicts; feature-distance metrics are necessary-not-sufficient"), not as an intentional experimental arm. All headline results use v2.

\subsection{7. Analyses}

\subsubsection{7.1 Multi-seed paired statistics}
The base+calibration protocol is run at three seeds (42, 43, 44) for k$\in$\{0,1,3\}, giving n = 15 (5 subjects $\times$ 3 seeds) per (method, k). We report per-(method,k) mean $\pm$ sd, paired $\Delta$ vs \texttt{scratch} with 95\% CIs and paired-t p-values (subject-and-seed-blocked). At n=5 a two-sided Wilcoxon bottoms out at p=0.0625, so subject-level tests cannot reach p<0.05 alone; multi-seed pooling and clip-level McNemar (\S{}7.2) carry significance instead. Selection rule: never best-test-epoch as primary; hyperparameters frozen on same-domain/validation data.

\subsubsection{7.2 Clip-level McNemar}
Because each held-out subject contributes \textasciitilde{}400--500 clips, we dump per-clip posteriors (predicted class, softmax confidence, correctness, logits) from each base+calibration checkpoint (\texttt{temp\_dump\_posteriors.py}) and run a paired \textbf{McNemar test} on prior-vs-\texttt{scratch} predictions, pooled over seeds. We report the discordant counts b (scratch-only correct) and c (prior-only correct), net = c $-$ b, and the McNemar p-value per (prior, k). No retraining --- this reuses the multi-seed checkpoints.

\subsubsection{7.3 Calibration (ECE)}
From the same posteriors we compute Expected Calibration Error (15-bin), before and after single-parameter temperature scaling (temperature fit per method$\times$k), and reliability diagrams. This quantifies how trustworthy each objective's confidence stream is --- the input the controller (\S{}8) consumes.

\subsubsection{7.4 Representation similarity (CKA)}
Layer-wise linear and RBF \textbf{Centered Kernel Alignment} between NTU and Xsens-v2 activations, per encoder, at the projection head and each of the three Transformer blocks (\texttt{temp\_cka\_analysis.py}). CKA is scale-invariant, removing the MMD scale confound. We also retain squared-MMD (median-heuristic and fixed-$\sigma$) as the necessary-not-sufficient companion. CKA is computed on \textbf{v2 only} (the swing arm is dead).

We also used the same layer-wise CKA to test whether the small/middle/large gap ordering across our five source$\rightarrow$target settings (R6e) is measurable rather than narrative: \texttt{temp\_cka\_analysis.py --multi-target} computes NTU-vs-target CKA for each of Xsens-v2, CZU-MHAD skeleton, CZU-MHAD IMU orientation quats, and UTD-MHAD skeleton (n$\approx$1000 clips/side, seed 0, inference-only). It does not confirm the ordering (Results R3) --- we report the attempt and the null result rather than omit it or force a fit.

\textbf{Raw, encoder-free domain gap.} CKA and the encoder-space MMD above are both computed inside a \textit{trained encoder's} feature space, confounding the measurement with the pretraining objective. To measure the datasets' distributional gap independent of any model, \texttt{temp\_raw\_domain\_gap.py} computes a hand-crafted, model-free feature per clip --- mean/std/var/skew/kurtosis over the 68 flattened LRQ channels (340-d; the same feature used for the CRC published-baseline reproductions, \S{}6/R6/R6d) --- for NTU and each of the four targets, then reports two metrics on pooled-standardized features (n=1000/side, n=861 for UTD, its full pool; seed 0): squared-MMD with an RBF kernel and median-heuristic bandwidth (fixing the fixed-$\sigma$ scale confound of the encoder-space MMD), and Frechet distance (Gaussian closed form on a 50-component PCA projection, for numerical stability of the covariance term at this sample size).

\subsubsection{7.5 Subject-count scaling (A2)}
To ask \textit{how long the prior stays useful as target data accumulates}, we sweep the number of fine-tuning subjects \texttt{N $\in$ \{0,1,2,3,4\}} used in the base phase (\texttt{--n-train-subjects}, taking the first N of the sorted non-held-out subjects; N=0 short-circuits to the pretrained init followed by k-shot calibration only; N=4 recovers the full protocol of \S{}5). 3 seeds (42/43/44), 5 LOSO folds per N per seed (n=15/cell), methods \{scratch, mae, supMAE, supLP120\}, k$\in$\{0,1,3\}. We report prior benefit vs scratch as a function of N (\texttt{temp\_a2\_run.sh} + \texttt{temp\_t1\_multiseed\_run.sh}, \texttt{temp\_analyze\_a2.py}). We additionally repeat this sweep (N$\in$\{0..3\}, seed 42, \{scratch, supLP120, supMAE\}) on the CZU dual-branch target (\S{}7's strong-target recognizer) to test whether the cold-start advantage generalizes beyond a representationally weak target.

\subsubsection{7.6 A fifth objective: supervised contrastive (SupCon)}
To complete the objective family named in the abstract, we add a fifth pretraining objective --- Khosla et al. (2020) SupCon, label-supervised contrastive learning (in-batch same-label positives) --- trained on NTU-120 identically to the other four, and re-run every analysis above (R1 linear probe and cross-domain LOSO-v2, R2 McNemar, R3 CKA, R4 ECE/AUC/convergence, R5 controller robustness, R6/R6b/R6c/R6d external datasets) with supcon included. Where an analysis depends on a fixed artifact shared across methods (the controller's reliability-ordered System Input assignment, \S{}8), adding a fifth method can silently perturb that artifact if its construction pools over all loaded methods rather than a fixed subset --- we discovered and document one such instance (Results R5) rather than either hiding it or discarding the extension.

\subsection{8. Control-reliability simulation (C6 / THMS pillar)}

To translate recognition accuracy and calibration into human-machine-systems reliability metrics, the models' posterior streams are evaluated through an abstract, event-driven Finite State Machine (FSM) that consumes the \textbf{actual recognizer's posterior stream} --- no physics engine (the scientific content is uncertainty$\rightarrow$reliability propagation; a simulator would add risk without signal). \texttt{temp\_controller\_sim.py}. States are defined purely by consequence, not by any physical semantics: a \textbf{Safety-Critical State} is one where a misclassification or false activation triggers immediate task failure; a \textbf{Routine State} allows a \textbf{Recoverable Error} at a time/effort penalty.

\begin{itemize}
  \item \textbf{Action-primitive assignment.} Seven abstract \textbf{Action Primitives} (System Inputs) make up the vocabulary for a \textbf{12-step Sequential Control Task}; none is itself a recorded gesture class --- none of the 22 recorded classes is defined as a System Input a priori. Seven of the 22 real gestures are instead \textit{relabeled} onto these primitives, purely by recognition reliability: gestures are ranked by mean recall pooled over \textit{all} methods at k=3, locomotion/whole-body-transition classes are excluded from primitive duty, and the two most reliably-recognized remaining gestures are assigned to the two \textbf{Safety-Critical States}, the rest to \textbf{Routine States}. The assignment is deliberately non-semantic (it tunes on aggregate reliability, favoring no single objective, rather than on which gesture "looks like" a given primitive) --- the task and its System Input assignment are a vehicle for studying how recognition and calibration propagate to task-level outcomes, not a claim that this dataset contains any specific real-world command gestures. Specific physical movements are therefore not named against states in what follows.
  \item \textbf{Sequential Control Task.} An ordered 12-step task containing three Safety-Critical steps (drawn from the two Safety-Critical States) and nine Routine steps, so recognition errors on safety-critical steps compound.
  \item \textbf{FSM + safety layer.} For each task step with an intended System Input c, a held-out clip of the mapped gesture is sampled and its recognizer posterior consumed. A \textbf{confidence-threshold reject} ($\tau$) re-prompts when max-softmax < $\tau$ (dwell/temporal-smoothing option requires agreement across repeated reads). A correct input advances the FSM; a \textbf{misclassification or false activation into a Safety-Critical State is an immediate task failure} (unsafe action); a Recoverable Error in a Routine State incurs a cancel-and-reissue cost. Excessive consecutive rejections abort the task.
  \item \textbf{Distractor / null stream.} Clips whose true class is \textit{not} a System Input drive the \textbf{false-activation rate} (fraction that pass $\tau$ and map to a System Input).
  \item \textbf{Metrics}, reported per (init $\times$ k) and swept over $\tau$: task success rate, mean time-to-completion (in cost units), rejection rate, corrective-inputs-per-task, and false-activation rate. Monte-Carlo (1000--3000 task runs/condition), fixed RNG seed.
\end{itemize}

\subsubsection{8.1 Robustness protocol --- locking the design knobs}
A controller has design parameters (System Input assignment, error-cost model, operating threshold) that could each be tuned to manufacture a result. Instead of freezing arbitrary values, and treating the \textbf{Monte Carlo randomized-assignment protocol as the primary evaluation protocol} (so findings are fundamental to the objective, not an artifact of interface design), we demonstrate the method \textbf{ordering is invariant} to all three, via three locks (\texttt{temp\_controller\_robust.py}; outputs \texttt{trained\_models/Phase3-controller/robust/}):

\begin{itemize}
  \item \textbf{Lock 1 --- randomized System Input assignment.} The 7-gesture$\rightarrow$primitive mapping is resampled at random 120 times; we report the \textit{distribution} of task success across assignments and the fraction of assignments in which each pairwise method ordering holds. This removes any dependence on a hand-picked System Input set.
  \item \textbf{Lock 2 --- critical-cost sweep + two outcome models.} In a single pass we compute both a \textbf{hard-safety} outcome (any Safety-Critical error $\rightarrow$ task failure, binary success) and a \textbf{soft-cost} outcome (Safety-Critical errors are recoverable but incur a penalty C\_crit), sweeping C\_crit $\in$ \{2, 5, 10, 20, 50, $\infty$\} ($\infty$ recovers hard-safety). This shows the ordering is independent of how catastrophic Safety-Critical errors are made.
  \item \textbf{Lock 3 --- iso-safety operating point.} Because real-world systems operate on strict false-activation budgets rather than a static accuracy threshold, we fix a false-activation \textit{budget} (1\%, 0.5\%) and find the smallest $\tau$ meeting it per method, comparing throughput at that fixed point --- a deployment-standard, tuning-free way to set thresholds. The full $\tau$-frontier (task-success vs false-activation) is also reported.
\end{itemize}

All reported controller claims come from this robustness protocol; no single fixed-assignment configuration is reported in isolation. The base cost weights (T\_exec = 1, T\_reject = 1, T\_correct = 3, max consecutive rejects = 5) are held fixed while the load-bearing knobs are swept as above.

\subsection{9. External-validity check --- CZU-MHAD skeleton$\rightarrow$skeleton}

As an independent public target we run the same objectives and LOSO protocol on CZU-MHAD's skeleton modality (1165 clips, 5 subjects, 22 actions; Kinect-v2 25-joint order mapped to NTU's via \texttt{temp\_czu\_parser.py} $\rightarrow$ \texttt{Data\_Processed/czu\_skeleton\_lrq/}). This is \textbf{skeleton$\rightarrow$skeleton} --- same modality, a much smaller gap than NTU$\rightarrow$Xsens IMU --- so it tests whether the objective orderings replicate on other data, not the IMU gap itself. CZU also ships 10 body-worn 6-axis inertial sensors (\texttt{sensor\_mat/}, no magnetometer); the true cross-modal replication via that IMU data is left as future work.

\subsection{10. Reproducibility notes}
Encoders, checkpoints (\texttt{encoder\_state\_dict}/\texttt{decoder\_state\_dict}/\texttt{head\_state\_dict} + metadata), and per-clip posteriors are retained under \texttt{trained\_models/}. Runs: v2 multi-seed \texttt{LOSO-fullTrainCalibrate-v2\{,-seed43,-seed44\}/}; Phase-1 analysis \texttt{Phase1-analysis/}; subject-scaling \texttt{A2-subjectScaling/}; controller \texttt{Phase3-controller/}; external validity \texttt{CZU-skeleton-LOSO/}. Data regeneration is deterministic given the parsers above.
